% =============================================================================
% SECTION 9: DISCUSSION AND FALSIFIABILITY
% =============================================================================
\section{Discussion and Falsifiability Summary}\label{sec:discussion}

The Universality Theorem asserts that foundation MLIPs in the architectural class $\cF$ share class-uniform geometric structure in their error manifolds. This section summarizes the implications, restates the pre-registered predictions, and identifies open mathematical problems.

\subsection{Implications for MLIP Evaluation}

The theorem reframes how foundation MLIPs should be evaluated. Current practice emphasizes aggregate error metrics (MAE, RMSE, F1) on held-out test sets. The Universality Theorem implies that these aggregates obscure the geometric structure of errors: two models with identical MAE may have error manifolds with different orientations (model-specific components) but the same intrinsic dimension (class-uniform property). The appropriate evaluation protocol should therefore include:
\begin{enumerate}
    \item \textbf{Residual PCA.} Decompose the cross-configuration residual matrix into principal directions and test for alignment across models (Prediction P2).
    \item \textbf{Fisher spectrum analysis.} Compute the empirical Fisher information matrix and test for Vandermonde decay at rate $\geq 1.5$ (clause iii).
    \item \textbf{Active learning curves.} Measure the excess-risk scaling with query count and test for $O(1/T)$ decay with constant within the Raj--Bach bound (Prediction P1).
\end{enumerate}

The MLIP Arena\cite{chiang2025mliparena} framework already implements elements of this protocol; the theorem provides the theoretical justification for expanding it to include geometric error-structure analysis.

\subsection{Implications for MLIP Improvement}

The three escape classes (Corollary~\ref{cor:escape}) provide a decision framework for allocating effort among improvement strategies. If the dominant error is class-uniform (paradigm-agnostic), then $\Delta$-learning on any model in $\cF$ yields comparable improvement, and the community should focus on developing better base embeddings. If the dominant error is model-specific, then fine-tuning is the appropriate strategy, and the community should focus on efficient fine-tuning protocols\cite{hanseroth2025finetuning}. If the error is structural inadequacy, then architectural innovation is required, and the refusal mode should be engaged.

The $\Delta$-learning bound (Proposition~\ref{prop:delta}) quantifies the trade-off: the error reduction factor is $O(d(\cF)^{-1/2})$, meaning that models with lower intrinsic dimension (more compressed error manifolds) benefit less from $\Delta$-learning because their residuals are already concentrated in a low-dimensional subspace. This counterintuitive result---that ``better'' models (in the sense of lower $d(\cF)$) are harder to improve via post-hoc correction---has practical implications for model development.

\subsection{Falsifiability Summary}

The theorem makes two pre-registered predictions with explicit falsification thresholds:

\textbf{Prediction P1 (clause iv).} The active-learning excess risk scales as $\epsilon(T) = C \cdot d(\cF) \log N / T$ with $C \in [C_{\text{RB}}/4, 4 C_{\text{RB}}]$. Falsification: $C$ outside this interval or scaling exponent $\neq -1 \pm 0.2$.

\textbf{Prediction P2 (clause vi).} Consecutive generations of a foundation MLIP family have top-5 error-direction cosine similarity $\geq 0.7$. Falsification: cosine similarity $< 0.5$ for any transition.

Both predictions are testable with publicly available models and benchmarks (Matbench Discovery, MLIP Arena). The falsifiability framework ties the abstract theorem to concrete empirical tests that any researcher can execute.

\subsection{Open Mathematical Problems}

Three open problems remain:

\textbf{Problem 1: Sharp class-uniform smoothness constant.} The constant $\kappa_3$ in (F3) is currently bounded by GAP-style smoothness analyses\cite{bigi2022smooth} but not by a published constant for equivariant-message-passing architectures. A sharp bound on $\kappa_3$ for MACE-style architectures would translate directly into a sharp bound on $\rho(\cF)$ via Equation~\eqref{eq:rho-bound}.

\textbf{Problem 2: Tightness of the Vandermonde lemma.} Is the constant $\rho(\cF) \leq C \kappa_1^{-1} \kappa_2 \kappa_3$ tight? That is, do there exist models $M \in \cF$ for which the decay rate achieves this upper bound, or can the bound be improved by a more refined analysis of the Beckermann conditioning parameter?

\textbf{Problem 3: Refusal region and incompleteness manifold.} The refusal region $\cX_N^{\text{refuse}}(M)$ is characterized geometrically via Federer reach, but its relationship to the Pozdnyakov--Ceriotti incompleteness manifold\cite{pozdnyakov2020incompleteness}---the set of configurations that cannot be distinguished by descriptors below a given body order---is not yet understood. Is the refusal region a tubular neighborhood of the incompleteness manifold? Answering this would connect the geometric refusal mode to the representational limitations documented by the incompleteness floor.

\subsection{Conclusion}

The Universality Theorem provides a geometric framework for understanding the error structure of foundation MLIPs. By proving that per-model error manifolds have class-uniform properties---bounded dimension, Vandermonde Fisher spectrum, and stable evolution---the theorem transforms the empirical observation of cross-model heterogeneity into a structured prediction about the geometry of machine-learned potential energy surfaces. The two pre-registered predictions and the three escape classes for error reduction provide actionable guidance for the development and evaluation of next-generation foundation MLIPs.
